% !PRINTED 282
\weberpage{282}
%<<CONT
\[
\begin{array}{ll}
\text{für } n=7, & m=24,\\
\textquotedblbase\quad n=23, & m=24,\\
\textquotedblbase\quad n=31, & m=96,\\
\textquotedblbase\quad n=47, & m=48,\\
\textquotedblbase\quad n=71, & m=72,\\
\textquotedblbase\quad n=15, & m=72.
\end{array}
\]

In den Fällen 23, 47, 71 ist also der Grad der Korrespondenz
gleich dem Grad der entsprechenden Modulargleichung, d. h. gleich
der Anzahl der betreffenden Transformationen; in den Fällen
$n=7, 31, 15$ ist der Grad das \gesperrt{Dreifache} des Grades der
Modulargleichungen.

Was haben diese überzähligen Punkte zu bedeuten? Sie
erklären sich dadurch, daß in diesen Fällen, in denen $n\equiv 0, 1$
$\pmod 3$ ist, in der Funktion $\Phi_{a,c,\partial}$ [§ 75, (5)], die nach der obigen
Gleichung (5) einer Konstanten gleich ist, die Exponenten $h+2k$
in allen Gliedern denselben Rest nach dem Modul 3 lassen, und
daß also die Gleichung $\Phi=0$ und ebenso die Gleichungen (2)
erfüllt bleiben, wenn $\xi, \eta, \zeta$ mit einer beliebigen dritten Einheitswurzel multipliziert werden. Es geben also je drei Punkte der
Korrespondenz dieselbe Transformation.

Wir hätten auch, wenn $x, y, z$ durch (1) bestimmt sind,
statt (4)

\weq{6}{\xi=f\left(\dfrac{c+\partial\omega}{a}\right),\ \eta=\left(\dfrac{2}{\partial}\right)f_2\left(\dfrac{c+\partial\omega}{a}\right),\ \zeta=\left(\dfrac{2}{a}\right)f_1\left(\dfrac{c+\partial\omega}{a}\right)}

und folglich für $A, B$:

\weq{7}{\begin{aligned} 2A &= x\xi+(-1)^{\frac{n+1}{8}}(y\zeta+z\xi),\\ B &= x\xi y\zeta+x\xi z\eta+(-1)^{\frac{n+1}{8}}y\eta z\zeta \end{aligned}} % [sic] Source prints z\xi in 2A; the symmetric pattern suggests z\eta.

setzen können. Die Grundkurve und der Grad der Korrespondenz
wären dann dieselben geblieben, aber die Koinzidenzen hätten
eine andere, und zwar einfachere, Bedeutung bekommen. In
beiden Fällen gehören die Koinzidenzen zu den singulären Werten
der Modulfunktionen, in denen $\omega$ eine imaginäre quadratische
Irrationalität ist, wie wir in der Folge noch genauer sehen
werden.

Im Falle $n\equiv 3 \pmod 8$ setzen wir
\[
x=f^2(\omega),\ y=f_1^2(\omega),\ z=f_2^2(\omega),
\]
\weq{8}{\xi=f^2\left(\dfrac{c+\partial\omega}{a}\right),\ \eta=f_1^2\left(\dfrac{c+\partial\omega}{a}\right),\ \varrho=f_2^2\left(\dfrac{c+\partial\omega}{a}\right),}
%CONT>>
